\section{Conclusion}
\label{sec:conclusion}

A pre-registered evaluation of a capacity-aware threshold detector
with $\tau_K = \{50: 0.20, 100: 0.05, 200: 0.02\}$ on a 2{,}700-row
sweep produces four hypothesis decisions: H1 supported (feasibility
region reached), H3 supported (firing perfectly monotone in K),
H2 rejected at one cell-window, H4 rejected (cap\_aware $=$ gated at
K=200 to four decimal places).

\textbf{The H1 positive.} The capacity-aware detector is the first
procedure in the Papers~7--12 sequence that simultaneously satisfies
all three pre-registered tolerances (worst K=200 $\geq -0.01$,
cell-mean K=50 $\geq -0.02$, cell-mean K=100 $\geq -0.02$). The
$\tau_{50} = 0.20$ never fires at K=50, so cap\_aware = lag-1 there
exactly; the $\tau_{200} = 0.02$ fires every window at K=200, so
cap\_aware = fixed there exactly; the K=100 intermediate produces a
$-0.004$ cell-mean deviation from lag-1.

\textbf{The H4 collapse.} cap\_aware and gated produce identical
cell means at K=200 ($0.2664$ each) because both strategies make the
same decision at every K=200 window: revert to Paper~4 fixed
weights. The per-window magnitude test contributes nothing beyond
what the unconditional gate already provides.

\textbf{The substantive interpretation.} Capacity-aware
thresholding reaches the feasibility region by approximating the
static gate at the two capacity extremes. The detector validates the
static gate as a feasible operating point but does not improve on
it in this grid. The deployable recipe is the static gate.

\textbf{Reproducibility.} Runner, analysis, pre-registration protocol,
and frozen 2{,}700-row results CSV are available at
\url{https://github.com/Harshavardhanmalla-Labs/research-papers/tree/main/paper12}.

\textbf{Future work.} (i) CUSUM~\cite{page1954cusum} detector with
pre-registered accumulator parameters; (ii) Bayesian online
change-point detector with pre-registered priors; (iii) per-K $\tau$
sensitivity sweep within the capacity-aware family; (iv) real fleet
validation with longitudinal exploitation outcome data; (v) closed-loop
adaptive procedures where the selection policy is also recalibrated.
